\section*{Appendix}

The appendixes are organized as follows:
\begin{itemize}
	\item Appendix~\ref{sec:related} provides a detailed discussion of the previous works most related to ours.
	\item Appendix~\ref{sec:hyperplane} provides all the details about  the \texttt{Find-HS} procedure which is employed by the \texttt{Discover-and-Cover} algorithm, including all the technical lemmas (and their proofs) related to such a procedure.
	\item Appendix~\ref{sec:app_algo} provides all the details about the \texttt{Discover-and-Cover} algorithm that are omitted from the main body of the paper, including the definitions of the parameters required by the algorithm and the proofs of all the results related to it.
	\item Appendix~\ref{sec:other_proofs} provides all the other proofs omitted from the main body of the paper.
\end{itemize}

\section{Related works}\label{sec:related}

In this section, we survey all the previous works that are most related to ours.
%
Among the works addressing principal-agent problems, we only discuss those focusing on learning aspects.
%
Notice that there are several works studying computational properties of principal-agent problems  which are \emph{not} concerned with learning aspects, such as, \emph{e.g.},~\citep{dutting2019simple,alon2021contracts,guruganesh2021contracts,dutting2021complexity,dutting2022combinatorial,castiglioni2022bayesian,castiglioni2023multi}.
%\begin{comment}
%\end{comment}




\paragraph{Learning in principal-agent problems}
%
The work that is most related to ours is~\citep{zhu2022sample}, which investigate a setting very similar to ours, though from an online learning perspective (see also our Section~\ref{sec:regret}).
%
% \citet{zhu2022sample} investigate a repeated principal-agent interaction, where the principal commits to contracts based on observed outcomes.
%
\citet{zhu2022sample} study general hidden-action principal-agent problem instances in which the principal faces multiple agent's types.
%
They provide a regret bound of the order of $\widetilde {\mathcal{O}}(\sqrt{m} \cdot T^{1 - 1/(2m+1)})$ when the principal is restricted to contracts in $[0,1]^m$, where $m$ is the number of outcomes.
%
They also show that their regret bound can be improved to $\widetilde {\mathcal{O}}( T^{1 - 1/(m+2)})$ by making additional structural assumptions on the problem instances, including the \emph{first-order stochastic dominance} (FOSD) condition.
%
Moreover, \citet{zhu2022sample} provide an (almost-matching) lower bound of $\Omega(T^{1-1/(m+2)})$ that holds even with a single agent's type, thus showing that the dependence on the number of outcomes $m$ at the exponent of $T$ is necessary in their setting.
%
Notice that the regret bound by~\citet{zhu2022sample} is ``close'' to a linear dependence on $T$, even when $m$ is very small.
%
In contrast, in Section~\ref{sec:regret} we show how our algorithm can be exploited to achieve a regret bound whose dependence on $T$ is of the order of $\widetilde{\mathcal{O}}(T^{4/5})$ (independent of $m$), when the number of agent's actions $n$ is constant.
%
% which is close to be linear even for small numbers of outcomes $m$. In contrast, we show that we can achieve a regret bound of $\mathcal{O}(T^{4/5})$ by exploiting a small agent's action space. 
%
Another work that is closely related to ours is the one by~\citet{ho2015adaptive}, who focus on designing a no-regret algorithm for a particular repeated principal-agent problem.
%
However, their approach relies heavily on stringent structural assumptions, such as the FOSD condition.
%
%on the well-known FOSD (First-Order Stochastic Dominance) assumption, which limits the generality of the setting.  
Finally, \cite{cohen2022learning} study a repeated principal-agent interaction with a \emph{risk-averse} agent, providing a no-regret algorithm that relies on the FOSD condition too.


\paragraph{Learning in Stackelberg games}
%
The learning problem faced in this paper is closely related to the problem of learning an optimal strategy to commit to in Stackelberg games, where the leader repeatedly interacts with a follower by observing the follower's best-response action at each iteration.
%
\citet{letchford2009learning} propose an algorithm for the problem requiring the leader to randomly sample a set of available strategies in order to determine agent's best-response regions.
%
The performances of such an algorithm depend on the volume of the smallest best-response region, and this considerably limits its generality.
%
\citet{Peng2019} build upon the work by~\citet{letchford2009learning} by proposing an algorithm with more robust performances, being them independent of the volume of the smallest best-response region.
%
The algorithm proposed in this paper borrows some ideas from that of~\citet{Peng2019}.
%
However, it requires considerable additional machinery to circumvent the challenge posed by the fact that the principal does \emph{not} observe agent's best-response actions, but only outcomes randomly sampled according to them.
%
%
% Our setting is closely related to repeated Stackelberg games, where a leader commits to strategies while observing the actions taken by a follower. In the context of learning the optimal leader's commitment, \cite{letchford2009learning} propose an algorithm that requires the principal to randomly sample a set of strategies in order to determine the agent's best-response regions. However, their algorithm's performance dependens on the volume of the smallest feasible region of the follower's action. In contrast,~\cite{Peng2019} introduce an algorithm that is independent of the volume of the smallest best-response region, offering more robust performance. We also notice that the algorithm we propose in the following can be seen as a generalization of the one developed by \cite{peng2019optimal}. However, our setting poses an additional challenge as we do not observe the action played by the agent, and thus, we design a suitably defined algorithms to overcame this issue.
%
%
Other related works in the Stackelberg setting are~\citep{bai2021sample}, which proposes a model where both the leader and the follower learn through repeated interaction, and~\citep{lauffer2022noregret}, which considers a scenario where the follower’s utility is unknown to the leader, but it can be linearly parametrized.
%
% They design a no-regret learning algorithm achieving a regret bound which is sublinear in the number of time steps, compared to the best policy in hindsight.
%
% In a more recent work~\cite{bai2021sample} propose a different setting where both the principal and the follower learn an equilibrium through repeated interactions. Finally, \cite{lauffer2022noregret} consider a setting in which the follower’s utility is unknown to the leader but it can be linearly parameterized. They design a no-regret learning algorithm achieving a regret bound which is sublinear in the number of time steps, compared to the best policy in hindsight. A remarkable example of repetaded security games is represented by repeated Stackelberg security games (\cite{balcan2009protecting},\cite{tambe2011security}), in which a defender (leader) strategically allocates its limited resources to protect targets against potential attackers (followers).
\paragraph{Assumptions relaxed compared to Stackelberg games}

In our work, we relax several limiting assumptions made in repeated Stackelber games (see, \emph{e.g.},~\citep{letchford2009learning,Peng2019}) to learn an optimal commitment. Specifically, in repeated Stackelberg games either the best response regions have at least a constant volume or they are empty. Thanks to Lemma~\ref{lem:solve_lpl}, we effectively address this limitation, showing that even when an optimal contract belongs to a zero-measured best-response region, we can still compute an approximately optimal solution. Furthermore, in Stackelberg games it is assumed that in cases where there are multiple best responses for the follower, any of them can be arbitrarily chosen. In contrast, we assume that the agent always breaks ties in favor of the leader as it is customary in the Stackelberg literature. Finally, our approach does not require the knowledge of the number of actions of the agent, differently from previously proposed algorithms.
